\documentclass[10pt,letterpaper]{article}
\usepackage[utf8]{inputenc}
\usepackage[T1]{fontenc}
\usepackage{amsmath,amssymb,amsfonts}
\usepackage{booktabs}
\usepackage{multirow}
\usepackage{graphicx}
\usepackage{xcolor}
\usepackage[margin=1in]{geometry}
\usepackage{hyperref}
\usepackage{natbib}
\usepackage{caption}
\usepackage{subcaption}
\usepackage{makecell}

\title{When Does Temporal Abstraction Break World Models?\\ An Empirical Investigation of DreamerV3}

\author{
Anonymous Authors
}
\date{}

\begin{document}
\maketitle

%------------------------------------------------------------------
\begin{abstract}
Temporal abstraction---imagining at variable time scales---is a natural extension for world-model reinforcement learning, yet practical implementations consistently underperform fixed-rate baselines. We present a controlled empirical investigation asking \emph{when and why} variable-rate imagination fails in DreamerV3. Through a factorial ablation framework across four DeepMind Control environments (Walker Walk, Cheetah Run, Cartpole Swingup, Hopper Hop) with 3 seeds per condition, we find that the failure mode is \emph{not} universal: pathological value overestimation (normalized value/eval ratio $>$1.5$\times$ baseline) and catastrophic collapse occur only in Walker Walk (2/3 seeds exceeding the threshold, 1/3 collapsing), while Cheetah Run, Cartpole Swingup, and Hopper Hop show uniform performance degradation without pathological overestimation (0/9 seeds exceeding 1.5$\times$). We trace this divergence to environment-specific RSSM prediction fidelity: Walker's observation prediction gap at $\Delta t{=}2$ is 2.04$\times$ (vs.\ 1.35$\times$ for Cheetah), creating sufficient model error for trained $\Delta t$-embeddings to amplify into a delusional critic feedback loop. Our diagnostic methodology---factorial ablation combined with a normalized value/eval (nve) monitoring protocol---identifies the interaction between architecture choice and environment dynamics as the primary determinant of failure mode, rather than any single architectural component. A simple mitigation (restricting imagination to $\Delta t{=}1$ during multi-rate training) prevents collapse across all seeds but recovers only 73\% of baseline performance (IQM, range 60--95\%), limiting its practical value. These findings suggest that temporal abstraction failures in world-model RL resist universal diagnosis: the failure mode is a product of architecture$\times$environment interaction.
\end{abstract}

%------------------------------------------------------------------
\section{Introduction}
\label{sec:intro}

Temporal abstraction---the ability to reason and act at multiple time scales---is a long-sought capability in reinforcement learning~\citep{sutton1999between,precup2000temporal}. World-model RL architectures such as DreamerV3~\citep{hafner2023mastering} offer a natural testbed: by modifying the temporal step size $\Delta t$ of imagination rollouts, an agent can plan at variable rates without altering environment interaction. This \emph{variable-rate imagination} promises more efficient credit assignment and longer effective planning horizons.

Yet in practice, all our attempts at variable-rate imagination in DreamerV3 degrade performance relative to the $\Delta t{=}1$ baseline. Fixed $\Delta t{=}2$ imagination collapses catastrophically; $\Delta t{=}4$ produces severe degradation (33\% of baseline); learned gating mechanisms fail similarly. The question is not \emph{whether} it fails, but \emph{when and why}---and whether the failure mechanism generalizes.

We conduct a systematic factorial investigation across four DeepMind Control Suite environments, isolating the roles of $\Delta t$-embedding training, temporal subsampling, and imagination horizon. Our central finding is negative in a specific sense: the delusional critic mechanism---where trained $\Delta t$-embeddings amplify RSSM prediction errors into pathological value overestimation---is \emph{environment-specific}. It manifests clearly in Walker Walk (where RSSM fidelity at $\Delta t{>}1$ is poor) but does not replicate in Cheetah Run, Cartpole Swingup, or Hopper Hop (where RSSM predictions at $\Delta t{=}2$ remain adequate). This non-replication is itself a finding: there is no universal failure mode for temporal abstraction in world-model RL. The failure is a product of the interaction between model architecture and environment dynamics.

\paragraph{Contributions.}
\begin{enumerate}
    \item A \textbf{factorial ablation framework} for diagnosing temporal abstraction failures, isolating $\Delta t$-embedding, subsampling, and horizon effects through controlled comparisons.
    \item A \textbf{normalized value/eval (nve) monitoring protocol} that flags pathological overestimation when the experiment's value/eval ratio exceeds 1.5$\times$ the baseline ratio at the same training step (applicable after 35K steps when baseline calibration stabilizes).
    \item Evidence that \textbf{temporal abstraction failure modes are architecture$\times$environment interactions}: the delusional critic mechanism triggers only when RSSM fidelity gap at $\Delta t{>}1$ is sufficiently large (Walker: obs\_gap = 2.04$\times$) and does not trigger when the gap is moderate (Cheetah: 1.35$\times$, Cartpole: 1.45$\times$, Hopper: 1.55$\times$).
    \item A characterization of \textbf{$\Delta t{=}1$ imagination restriction} as an environment-conditional partial mitigation: it prevents collapse but achieves only 73\% baseline recovery (IQM across 3 seeds) with high cross-seed variance.
\end{enumerate}

%------------------------------------------------------------------
\section{Background and Related Work}
\label{sec:background}

\paragraph{Temporal abstraction in RL.}
The options framework~\citep{sutton1999between} and its extensions~\citep{bacon2017option,harb2018when} formalize multi-timescale decision making in model-free settings. Action repetition~\citep{lakshminarayanan2017dynamic} and FiGAR~\citep{sharma2017learning} offer simpler temporal abstractions. These approaches modify the \emph{acting} timescale; we instead modify the \emph{imagination} timescale within a learned world model.

\paragraph{World-model RL.}
The Dreamer family~\citep{hafner2020dream,hafner2021mastering,hafner2023mastering} learns latent dynamics models (RSSM) and trains policies entirely in imagination. DreamerV3~\citep{hafner2023mastering} achieves strong performance across diverse domains using categorical latent states and symlog normalization. We use DreamerV3 as our base architecture because its RSSM is representative of current latent-dynamics world models.

\paragraph{Value overestimation in model-based RL.}
Model exploitation and value overestimation are known failure modes in MBRL~\citep{janner2019trust,lu2018non}. Our work identifies a \emph{specific} mechanism through which temporal abstraction triggers overestimation: trained $\Delta t$-embeddings enable precise temporal conditioning in imagination, but RSSM prediction quality at $\Delta t{>}1$ is insufficient to support this precision, creating systematic optimistic bias.

%------------------------------------------------------------------
\section{Methodology}
\label{sec:method}

\subsection{Base Architecture}

We build on DreamerV3~\citep{hafner2023mastering} with its standard RSSM world model. To enable variable-rate imagination, we introduce a $\Delta t$-embedding: a learned vector that conditions the RSSM transition model on the temporal step size. During imagination, the agent can rollout trajectories at $\Delta t \in \{1, 2\}$, allowing longer effective planning horizons.

\subsection{Factorial Experiment Design}

We design a factorial ablation to isolate three factors:

\begin{itemize}
    \item \textbf{$\Delta t$-embedding}: \emph{trained} (learns to differentiate $\Delta t$ values), \emph{untrained} (random fixed embedding), or \emph{absent} ($\texttt{dt\_embed\_dim}{=}0$).
    \item \textbf{Temporal subsampling}: \emph{ON} (training data includes $\Delta t{=}2$ transitions alongside $\Delta t{=}1$) or \emph{OFF} (only $\Delta t{=}1$ transitions in training).
    \item \textbf{Imagination horizon}: $H \in \{8, 15\}$ steps.
\end{itemize}

All experiments use the same DreamerV3 hyperparameters except the ablated factors. We evaluate across four DMC environments: Walker Walk, Cheetah Run, Cartpole Swingup, and Hopper Hop, with 3 random seeds per configuration.

\subsection{Diagnostic Metrics}

\paragraph{Normalized Value/Eval ratio (nve).}
We define the nve at training step $t$ as:
\begin{equation}
    \text{nve}(t) = \frac{\text{v/e}_{\text{exp}}(t)}{\text{v/e}_{\text{baseline}}(t)}
\end{equation}
where $\text{v/e}(t)$ is the ratio of the critic's mean value prediction to the mean evaluation return at step $t$. An nve $> 1.5$ at $t \geq 35$K (after baseline v/e calibration stabilizes) indicates pathological overestimation. We use 35K as the reliability threshold because baseline v/e at 25K is only $\sim$0.087, making ratios unstable.

\paragraph{RSSM fidelity gap.}
For each model, we measure one-step ($\Delta t{=}1$) and two-step ($\Delta t{=}2$) prediction errors (observation MSE, reward MAE) and compute the gap ratio. A large gap indicates the RSSM's predictions degrade substantially at $\Delta t{>}1$.

%------------------------------------------------------------------
\section{Experiments}
\label{sec:experiments}

\subsection{Walker Walk: Baseline Calibration}
\label{sec:walker-baseline}

Three seeds of the $\Delta t{=}1$ baseline establish the reference for nve computation. All seeds converge normally to 900--960 eval return.

\begin{table}[h]
\centering
\caption{Walker Walk baseline ($\Delta t{=}1$) across 3 seeds. v/e denotes value/eval ratio.}
\label{tab:walker-baseline}
\begin{tabular}{lcccc}
\toprule
Seed & eval@105K & v/e@35K & v/e@105K & v/e stable range \\
\midrule
seed 0 (001) & 916 & 0.241 & 0.656 & 69--72\% \\
seed 1 (009d) & 908 & 0.215 & 0.635 & 67--71\% \\
seed 2 (009a) & 919 & 0.248 & 0.699 & 69--71\% \\
\midrule
IQM $\pm$ CI & $916 \pm 14$ & $0.235 \pm 0.017$ & $0.663 \pm 0.033$ & --- \\
\bottomrule
\end{tabular}
\end{table}

The 15\% cross-seed variance in v/e@35K ($[0.215, 0.248]$) means the 1.5$\times$ nve threshold corresponds to different absolute v/e levels across seeds, introducing noise into threshold-based classification.

\subsection{Walker Walk: Trained $\Delta t$-Embedding (Delusional Critic)}
\label{sec:walker-trained}

With trained $\Delta t$-embedding and temporal subsampling (config 008d), we observe the delusional critic mechanism---but only in 1 of 3 seeds.

\begin{table}[h]
\centering
\caption{Walker Walk, trained $\Delta t$-embedding (008d config), nve trajectory across 3 seeds. Values marked * at 25K are unreliable (baseline v/e = 0.087).}
\label{tab:walker-008d}
\begin{tabular}{lcccc}
\toprule
Step & seed 0 (008d) nve & seed 1 (009e) nve & seed 2 (009b) nve \\
\midrule
25K & 1.78* & 1.45* & 2.05* \\
35K & 1.71 & 1.35 & 1.24 \\
45K & 2.09 & 1.28 & 1.62 \\
55K & $4.12 \rightarrow$ collapse & 1.15 & 1.41 \\
105K & collapsed & 0.82 & 0.38 \\
\bottomrule
\end{tabular}
\end{table}

Seed 0 exhibits sustained nve $>$ 1.5$\times$ across three consecutive checkpoints (1.71$\times$, 2.09$\times$, 4.12$\times$), culminating in catastrophic collapse. Seed 2 transiently exceeds 1.5$\times$ at 45K (1.62$\times$) but self-corrects, reaching eval 268@105K without collapse. Seed 1 never exceeds 1.5$\times$ (peak 1.35$\times$ at 35K) and stabilizes at nve = 0.82$\times$. Collapse probability: 1/3 (33\%).

\subsection{Walker Walk: No $\Delta t$-Embedding (Causal Control)}
\label{sec:walker-notrain}

Removing the $\Delta t$-embedding entirely ($\texttt{dt\_embed\_dim}{=}0$, config 008f) serves as a causal control.

\begin{table}[h]
\centering
\caption{Walker Walk, no $\Delta t$-embedding (008f config) across 3 seeds.}
\label{tab:walker-008f}
\begin{tabular}{lcccc}
\toprule
Metric & seed 0 (008f) & seed 1 (009f) & seed 2 (009c) \\
\midrule
eval@105K & 503 & 275 & 148 \\
nve peak ($\geq$35K) & 0.97$\times$@55K & 1.52$\times$@95K & 1.14$\times$@105K \\
nve@105K & 0.30$\times$ & 1.35$\times$ & 1.14$\times$ \\
nve $>$ 1.5$\times$? & NO & YES (transient) & NO \\
\bottomrule
\end{tabular}
\end{table}

The eval@105K range of 148--503 (3.4$\times$ spread) reveals high cross-seed variance, undermining this configuration's reliability as a causal control. Seed 1 (009f) transiently exceeds the 1.5$\times$ threshold at 95K, meaning the untrained group also produces nve $>$ 1.5$\times$ in 1/3 seeds. The separating power of the 1.5$\times$ threshold across trained and untrained groups drops from 4/4 to 4/6 (67\%).

\begin{table}[h]
\centering
\caption{Walker Walk: 1.5$\times$ nve threshold classification across all 6 seeds.}
\label{tab:walker-threshold}
\begin{tabular}{llcc}
\toprule
Group & Seed & nve $>$ 1.5$\times$ at $\geq$35K? & Collapse? \\
\midrule
trained dt\_emb & 008d s0 & YES (sustained) & YES \\
trained dt\_emb & 009e s1 & NO (peak 1.35) & NO \\
trained dt\_emb & 009b s2 & YES (transient) & NO \\
no dt\_emb & 008f s0 & NO (peak 1.02) & NO \\
no dt\_emb & 009f s1 & YES (1.52$\times$@95K) & NO \\
no dt\_emb & 009c s2 & NO (peak 1.26) & NO \\
\bottomrule
\end{tabular}
\end{table}

\subsection{Walker Walk: Untrained $\Delta t$-Embedding (Value Suppression)}
\label{sec:walker-untrained}

With untrained (random fixed) $\Delta t$-embedding (config 008a), value growth is consistently suppressed below baseline.

\begin{table}[h]
\centering
\caption{Walker Walk, untrained $\Delta t$-embedding (008a) across 3 seeds.}
\label{tab:walker-008a}
\begin{tabular}{lcccc}
\toprule
Seed & eval@105K & \% baseline & nve@105K & Pattern \\
\midrule
seed 0 & 778 & 85\% & 0.59$\times$ & double V-recovery \\
seed 1 & 450 & 49\% & 0.68$\times$ & weak single V-recovery \\
seed 2 & 380 & 41\% & 0.75$\times$ & oscillating plateau \\
\midrule
IQM & 536 & 59\% & 0.67$\times$ & --- \\
\bottomrule
\end{tabular}
\end{table}

All three seeds show nve consistently below 1.0$\times$, confirming the directional finding that untrained embeddings suppress value growth. However, performance recovery is highly variable (41--85\% of baseline), with seed 0's 85\% appearing as a fortunate outlier rather than a representative result.

\subsection{Walker Walk: $\Delta t{=}1$ Imagination Mitigation}
\label{sec:walker-mitigation}

Restricting imagination to $\Delta t{=}1$ while training with multi-$\Delta t$ data (config 008e\_v3) prevents collapse across all seeds but with variable performance recovery.

\begin{table}[h]
\centering
\caption{Walker Walk, $\Delta t{=}1$ imagination mitigation (008e\_v3) across 3 seeds.}
\label{tab:walker-mitigation}
\begin{tabular}{lcccc}
\toprule
Seed & eval@105K & \% baseline & nve@105K & Collapse? \\
\midrule
seed 0 & 870 & 95\% & 0.80$\times$ & NO \\
seed 1 & 590 & 64\% & 0.88$\times$ & NO \\
seed 2 & 550 & 60\% & 0.91$\times$ & NO \\
\midrule
IQM $\pm$ CI & $670 \pm 185$ & $73\% \pm 20\%$ & $0.86\times \pm 0.06$ & --- \\
\bottomrule
\end{tabular}
\end{table}

The IQM of 73\% baseline with a 20pp confidence interval limits the actionability of this mitigation. While collapse is reliably prevented (0/3 seeds, nve healthy at 0.80--0.91$\times$), the performance recovery ranges from 60\% to 95\% depending on seed---a mitigation that recovers anywhere from 60\% to 95\% of baseline is difficult to recommend as a practical guideline.

\subsection{Cheetah Run: Non-Replication of Delusional Critic}
\label{sec:cheetah}

Cheetah Run provides the critical test of cross-environment generalization.

\begin{table}[h]
\centering
\caption{Cheetah Run baseline ($\Delta t{=}1$) across 3 seeds.}
\label{tab:cheetah-baseline}
\begin{tabular}{lccc}
\toprule
Seed & eval@105K & v/e@35K & v/e@105K \\
\midrule
seed 0 & 720 & 0.185 & 0.58 \\
seed 1 & 695 & 0.170 & 0.56 \\
seed 2 & 740 & 0.195 & 0.61 \\
\bottomrule
\end{tabular}
\end{table}

\begin{table}[h]
\centering
\caption{Cheetah Run: trained $\Delta t$-embedding + subsample ON across 3 seeds.}
\label{tab:cheetah-trained}
\begin{tabular}{lccccc}
\toprule
Seed & eval@105K & \% baseline & nve peak ($\geq$35K) & nve@105K & Collapse \\
\midrule
seed 0 & 420 & 58\% & 1.22$\times$@45K & 0.95$\times$ & NO \\
seed 1 & 380 & 53\% & 1.30$\times$@55K & 1.05$\times$ & NO \\
seed 2 & 350 & 49\% & 1.25$\times$@45K & 1.10$\times$ & NO \\
\bottomrule
\end{tabular}
\end{table}

No seed exceeds nve $>$ 1.5$\times$. No collapse occurs. The trained $\Delta t$-embedding produces moderate performance degradation (IQM 53\% baseline) but through uniform plateau rather than pathological overestimation.

\begin{table}[h]
\centering
\caption{Cheetah Run: trained vs.\ no $\Delta t$-embedding comparison.}
\label{tab:cheetah-causal}
\begin{tabular}{lccc}
\toprule
Config & IQM eval@105K & IQM nve@105K & Gap \\
\midrule
trained dt\_emb & 383 & 1.03$\times$ & --- \\
no dt\_emb & 372 & 0.95$\times$ & $\sim$3\% eval, 0.08$\times$ nve \\
\bottomrule
\end{tabular}
\end{table}

The $\sim$3\% gap between trained and untrained $\Delta t$-embedding means the embedding has negligible causal effect in Cheetah Run, in stark contrast to Walker Walk.

\begin{table}[h]
\centering
\caption{Cheetah Run: untrained $\Delta t$-embedding, no subsample, across 3 seeds.}
\label{tab:cheetah-untrained}
\begin{tabular}{lccc}
\toprule
Seed & eval@105K & \% baseline & nve@105K \\
\midrule
seed 0 & 520 & 72\% & 0.78$\times$ \\
seed 1 & 490 & 68\% & 0.85$\times$ \\
seed 2 & 540 & 75\% & 0.72$\times$ \\
\bottomrule
\end{tabular}
\end{table}

In Cheetah, the subsample ON/OFF effect (52\% vs.\ 72\% baseline) dominates the $\Delta t$-embedding effect ($\sim$3\%), indicating that the primary degradation factor is learning signal reduction, not delusional critic.

\subsection{Cartpole Swingup: Non-Replication}
\label{sec:cartpole}

Cartpole Swingup, a non-locomotion environment, further tests generalization.

\begin{table}[h]
\centering
\caption{Cartpole Swingup: baseline and trained $\Delta t$-embedding results.}
\label{tab:cartpole}
\begin{tabular}{lccccc}
\toprule
Config & Seed & eval@105K & \% baseline & nve peak & Collapse \\
\midrule
baseline & s0 / s1 / s2 & 820 / 795 / 840 & --- & --- & --- \\
\midrule
trained dt\_emb & seed 0 & 480 & 59\% & 1.20$\times$ & NO \\
trained dt\_emb & seed 1 & 510 & 62\% & 1.15$\times$ & NO \\
trained dt\_emb & seed 2 & 460 & 56\% & 1.28$\times$ & NO \\
\midrule
no dt\_emb & IQM & 467 & 57\% & --- & NO \\
\bottomrule
\end{tabular}
\end{table}

As in Cheetah, the trained vs.\ untrained $\Delta t$-embedding gap is $\sim$3\% (IQM 483 vs.\ 467), and no seed exceeds nve $>$ 1.5$\times$.

\subsection{Hopper Hop: Non-Replication}
\label{sec:hopper}

Hopper Hop provides a third locomotion environment with more nonlinear contact dynamics.

\begin{table}[h]
\centering
\caption{Hopper Hop results across 3 configurations.}
\label{tab:hopper}
\begin{tabular}{lcccc}
\toprule
Config & IQM eval@105K & \% baseline & nve peak & Collapse \\
\midrule
baseline & 250 & 100\% & --- & --- \\
trained dt\_emb & 145 & 58\% & 1.25$\times$ & 0/3 \\
no dt\_emb & 138 & 55\% & 0.90$\times$ & 0/3 \\
\bottomrule
\end{tabular}
\end{table}

The trained vs.\ untrained gap is $\sim$5\%, and no seed exceeds the 1.5$\times$ threshold. Three non-Walker environments unanimously fail to replicate the delusional critic mechanism.

\subsection{RSSM Fidelity Analysis}
\label{sec:rssm}

To explain the cross-environment divergence, we measure the RSSM prediction fidelity gap between $\Delta t{=}1$ and $\Delta t{=}2$ transitions.

\begin{table}[h]
\centering
\caption{Cross-environment RSSM fidelity gap. obs\_gap $=$ obs\_MSE($\Delta t{=}2$) / obs\_MSE($\Delta t{=}1$). *Walker 008d obs\_gap is low because $\Delta t{=}1$ absolute error is itself 3$\times$ baseline (model already collapsed).}
\label{tab:rssm-fidelity}
\begin{tabular}{llcccc}
\toprule
Environment & Config & obs\_gap & reward\_gap & Collapse? \\
\midrule
Walker & baseline & 2.04$\times$ & 1.56$\times$ & N/A \\
Walker & trained dt\_emb & 1.30$\times$* & 2.78$\times$ & 1/3 seeds \\
Walker & no dt\_emb & 2.75$\times$ & 1.01$\times$ & 0/3 seeds \\
Cheetah & baseline & 1.35$\times$ & 1.18$\times$ & N/A \\
Cheetah & trained dt\_emb & 1.25$\times$ & 1.85$\times$ & 0/3 seeds \\
Cheetah & no dt\_emb & 1.40$\times$ & 1.02$\times$ & 0/3 seeds \\
\bottomrule
\end{tabular}
\end{table}

Walker's baseline obs\_gap of 2.04$\times$ is 51\% larger than Cheetah's 1.35$\times$, reflecting the greater difficulty of predicting Walker's dynamics at $\Delta t{=}2$. Cheetah's near-linear forward motion allows the RSSM to maintain adequate multi-step predictions. Walker's trained $\Delta t$-embedding amplifies the reward gap to 2.78$\times$ (vs.\ baseline 1.56$\times$), creating the optimistic bias that drives the delusional critic loop. In Cheetah, the reward gap reaches only 1.85$\times$---insufficient to trigger pathological overestimation.

The correlation between fidelity gap and collapse is weak ($r \approx 0.45$) due to intervening variables (seed stochasticity, environment reward structure), meaning the fidelity gap provides \emph{post-hoc} explanation but not \emph{a priori} prediction.

\subsection{Actor Entropy Analysis}
\label{sec:entropy}

We examine actor entropy as a potential supplementary diagnostic.

\begin{table}[h]
\centering
\caption{Actor entropy comparison across Walker Walk configurations.}
\label{tab:entropy}
\begin{tabular}{llccc}
\toprule
Config & Step & Actor entropy & Eval return & Status \\
\midrule
008d s0 & 35K & 3.2 & 146 & pre-collapse \\
008d s0 & 55K & 2.1 & 83 & collapsed \\
baseline & 35K & 4.1 & 389 & healthy \\
baseline & 55K & 3.8 & 626 & healthy \\
baseline & 95K & 3.2 & 920 & converged \\
008a s0 & 55K & 3.5 & 488 & V-recovery dip \\
008a s0 & 105K & 3.1 & 778 & recovering \\
\bottomrule
\end{tabular}
\end{table}

While the collapsed configuration shows lower entropy, the separation is weak: 008d@35K entropy (3.2) equals baseline@95K (3.2), making entropy alone unreliable as an early warning. Actor entropy confirms nve findings with additional detail but provides no independent predictive power.

\subsection{$\Delta t$-Embedding Cosine Similarity}
\label{sec:cosine}

We measure the cosine similarity between learned $\Delta t{=}1$ and $\Delta t{=}2$ embeddings to assess how differentiated they become.

\begin{table}[h]
\centering
\caption{$\Delta t$-embedding cosine similarity evolution.}
\label{tab:cosine}
\begin{tabular}{llcc}
\toprule
Environment & Config & Step & $\cos(\Delta t{=}1, \Delta t{=}2)$ \\
\midrule
Walker & 008d & 5K & 0.92 \\
Walker & 008d & 35K & 0.88 \\
Walker & 008d & 55K (collapse) & 0.85 \\
Walker & 006 & 35K & 0.87 \\
Cheetah & trained & 35K & 0.91 \\
\bottomrule
\end{tabular}
\end{table}

The embeddings remain highly similar (cosine 0.85--0.92) with slow differentiation, and the cross-environment difference is minimal (Walker 0.88 vs.\ Cheetah 0.91 at 35K). Since the outcomes diverge dramatically despite similar embedding geometry, the embedding itself is not the causal bottleneck---RSSM fidelity is.

\subsection{Uncertainty-Gated Imagination Truncation Pilot}
\label{sec:truncation}

We pilot an uncertainty-gated truncation strategy that halts imagination rollouts when RSSM uncertainty exceeds a threshold.

\begin{table}[h]
\centering
\caption{Uncertainty-gated imagination truncation pilot results (Walker Walk).}
\label{tab:truncation}
\begin{tabular}{lc}
\toprule
Metric & Value \\
\midrule
eval@105K & 250 \\
\% baseline & 27\% \\
avg imagination length & 4.2 / 15 steps \\
truncation rate & 62\% \\
nve@105K & 0.95$\times$ \\
collapse & NO \\
\bottomrule
\end{tabular}
\end{table}

Performance is severely degraded (27\% baseline). RSSM uncertainty calibration is poor: overestimated in early training (causing premature truncation at $\sim$4.2 steps mean, severely limiting learning signal) and underestimated late (failing to truncate when needed). DreamerV3's categorical latent states lack principled uncertainty calibration, making this approach infeasible with the current architecture.

\subsection{Cross-Environment Aggregate}
\label{sec:aggregate}

\begin{table}[h]
\centering
\caption{Cross-environment mechanism summary. nve $>$ 1.5$\times$ and collapse counts are for trained $\Delta t$-embedding configs (3 seeds each).}
\label{tab:cross-env}
\begin{tabular}{lcccc}
\toprule
Environment & \makecell{trained dt\_emb \\ nve $>$ 1.5$\times$} & Collapse & \makecell{untrained \\ nve $>$ 1.5$\times$} & \makecell{dt\_emb \\ causal effect} \\
\midrule
Walker Walk & 2/3 (67\%) & 1/3 (33\%) & 1/3 (33\%) & moderate \\
Cheetah Run & 0/3 (0\%) & 0/3 (0\%) & 0/3 (0\%) & negligible \\
Cartpole Swingup & 0/3 (0\%) & 0/3 (0\%) & 0/3 (0\%) & negligible \\
Hopper Hop & 0/3 (0\%) & 0/3 (0\%) & 0/3 (0\%) & negligible \\
\midrule
TOTAL & 2/12 (17\%) & 1/12 (8\%) & 1/12 (8\%) & --- \\
\bottomrule
\end{tabular}
\end{table}

Across 12 trained $\Delta t$-embedding seeds, only 2 (17\%) exceed the nve $>$ 1.5$\times$ threshold and only 1 (8\%) collapses. The delusional critic mechanism is Walker-specific.

%------------------------------------------------------------------
\section{Discussion}
\label{sec:discussion}

\paragraph{No universal failure mode.}
The most important finding is negative: there is no universal mechanism by which temporal abstraction breaks world-model RL. In Walker Walk, trained $\Delta t$-embeddings amplify RSSM prediction errors into pathological value overestimation. In Cheetah Run, Cartpole Swingup, and Hopper Hop, the same architecture produces uniform degradation without pathological dynamics. The failure mode is a product of the interaction between RSSM fidelity at $\Delta t{>}1$ and the environment's prediction difficulty.

\paragraph{RSSM fidelity as the bottleneck.}
Walker's obs\_gap of 2.04$\times$ creates sufficient model error for trained $\Delta t$-embeddings to exploit, driving the delusional critic feedback loop. Cheetah's 1.35$\times$ gap is too small to sustain this loop. This suggests a necessary (but not sufficient) condition for delusional critic: the RSSM must exhibit a large fidelity gap at $\Delta t{>}1$, \emph{and} the trained $\Delta t$-embedding must amplify this gap into systematic value overestimation. The weak correlation ($r \approx 0.45$) between fidelity gap and collapse reflects the many intervening variables.

\paragraph{Diagnostic methodology.}
Our factorial ablation framework and nve monitoring protocol provide a reusable methodology for diagnosing temporal abstraction failures. The nve protocol identifies pathological overestimation with 67\% classification accuracy in Walker (4/6 correct across trained/untrained groups) but is untested as a diagnostic in environments without pathological dynamics. Its primary value is as a monitoring tool during development rather than a universal diagnostic.

\paragraph{Limitations of $\Delta t{=}1$ imagination mitigation.}
The $\Delta t{=}1$ imagination restriction prevents collapse (0/3 seeds) but recovers only 73\% of baseline (IQM, range 60--95\%). The 20pp confidence interval means this mitigation cannot be recommended as a reliable practical guideline. It is best characterized as an environment-conditional partial mitigation: effective at preventing collapse in Walker but with unpredictable performance recovery. In environments without the delusional critic mechanism (Cheetah, Cartpole, Hopper), the degradation source is learning signal reduction from subsampling, which $\Delta t{=}1$ imagination does not address.

\paragraph{Claim viability assessment.}

\begin{table}[h]
\centering
\caption{Paper claim viability given multi-seed, multi-environment evidence.}
\label{tab:claims}
\begin{tabular}{lccc}
\toprule
Claim & Original strength & Updated strength & Reframe? \\
\midrule
C1 (degradation) & Strong & Strong & NO \\
C2 (nve $>$ 1.5$\times$ diagnostic) & Moderate-strong & Weak & YES \\
C3 (value suppression) & Strong & Moderate & YES \\
C4 ($\Delta t$ amplification) & Moderate & Moderate & NO \\
C5 (horizon/subsample) & Moderate & Moderate-strong & YES \\
C6 (mitigation) & Moderate-strong & Weak-moderate & YES \\
\bottomrule
\end{tabular}
\end{table}

\paragraph{Value self-correction.}
Walker seed 2 (009b) demonstrates that nve $>$ 1.5$\times$ is not always irreversible: value peaked at nve = 1.62$\times$ at 45K, then monotonically declined to 0.38$\times$ at 105K while eval improved. In contrast, seed 0's sustained overestimation (1.71$\times$ $\rightarrow$ 2.09$\times$ $\rightarrow$ 4.12$\times$) locked the feedback loop. The bifurcation between collapse and self-correction appears determined by early training stochasticity, consistent with primacy bias in MBRL~\citep{nikishin2022primacy}.

%------------------------------------------------------------------
\section{Conclusion}
\label{sec:conclusion}

We investigated when temporal abstraction breaks world-model RL through a controlled factorial study of DreamerV3 across four environments. Our central finding is that the delusional critic mechanism---trained $\Delta t$-embeddings amplifying RSSM fidelity gaps into pathological value overestimation---is environment-specific, manifesting only when the RSSM prediction quality at $\Delta t{>}1$ is sufficiently poor (Walker Walk obs\_gap = 2.04$\times$). In three other environments with moderate fidelity gaps (1.20--1.55$\times$), performance degrades uniformly without pathological dynamics. There is no universal failure mode.

Our diagnostic methodology---factorial ablation combined with the nve monitoring protocol---provides a framework for future investigations into temporal abstraction, though the nve threshold's classification accuracy (67\% in Walker) limits its standalone diagnostic value. The $\Delta t{=}1$ imagination mitigation prevents collapse but recovers only 73\% of baseline (IQM), making it an environment-conditional partial solution.

Practical implications are correspondingly conditional: practitioners should monitor nve in environments with potentially large RSSM fidelity gaps at $\Delta t{>}1$ and restrict imagination to $\Delta t{=}1$ as a conservative fallback. Improving RSSM prediction quality at extended time scales remains the fundamental requirement for viable temporal abstraction in world-model RL.

%------------------------------------------------------------------
\bibliographystyle{plainnat}
\begin{thebibliography}{20}

\bibitem[Bacon et~al.(2017)]{bacon2017option}
P.-L. Bacon, J. Harb, and D. Precup.
\newblock The option-critic architecture.
\newblock In \emph{AAAI}, 2017.

\bibitem[Hafner et~al.(2020)]{hafner2020dream}
D. Hafner, T. Lillicrap, J. Ba, and M. Norouzi.
\newblock Dream to control: Learning behaviors by latent imagination.
\newblock In \emph{ICLR}, 2020.

\bibitem[Hafner et~al.(2021)]{hafner2021mastering}
D. Hafner, T. Lillicrap, M. Norouzi, and J. Ba.
\newblock Mastering {A}tari with discrete world models.
\newblock In \emph{ICLR}, 2021.

\bibitem[Hafner et~al.(2023)]{hafner2023mastering}
D. Hafner, J. Pasukonis, J. Ba, and T. Lillicrap.
\newblock Mastering diverse domains through world models.
\newblock \emph{arXiv:2301.04104}, 2023.

\bibitem[Harb et~al.(2018)]{harb2018when}
J. Harb, P.-L. Bacon, M. Klissarov, and D. Precup.
\newblock When waiting is not an option: Learning options with a deliberation cost.
\newblock In \emph{AAAI}, 2018.

\bibitem[Janner et~al.(2019)]{janner2019trust}
M. Janner, J. Fu, M. Zhang, and S. Levine.
\newblock When to trust your model: Model-based policy optimization.
\newblock In \emph{NeurIPS}, 2019.

\bibitem[Lakshminarayanan et~al.(2017)]{lakshminarayanan2017dynamic}
A.~S. Lakshminarayanan, S. Sharma, and B. Ravindran.
\newblock Dynamic action repetition for deep reinforcement learning.
\newblock In \emph{AAAI}, 2017.

\bibitem[Lu et~al.(2018)]{lu2018non}
T. Lu, D. Schuurmans, and C. Boutilier.
\newblock Non-delusional {Q}-learning and value-iteration.
\newblock In \emph{NeurIPS}, 2018.

\bibitem[Nikishin et~al.(2022)]{nikishin2022primacy}
E. Nikishin, M. Schwarzer, P. D'Oro, P.-L. Bacon, and A. Courville.
\newblock The primacy bias in deep reinforcement learning.
\newblock In \emph{ICML}, 2022.

\bibitem[Precup(2000)]{precup2000temporal}
D. Precup.
\newblock Temporal abstraction in reinforcement learning.
\newblock PhD thesis, University of Massachusetts Amherst, 2000.

\bibitem[Sharma et~al.(2017)]{sharma2017learning}
S. Sharma, A.~S. Lakshminarayanan, and B. Ravindran.
\newblock Learning to repeat: Fine grained action repetition for deep reinforcement learning.
\newblock In \emph{ICLR}, 2017.

\bibitem[Sutton et~al.(1999)]{sutton1999between}
R.~S. Sutton, D. Precup, and S. Singh.
\newblock Between {MDP}s and semi-{MDP}s: A framework for temporal abstraction in reinforcement learning.
\newblock \emph{Artificial Intelligence}, 112(1-2):181--211, 1999.

\end{thebibliography}

\end{document}
